\documentclass[../../main.tex]{subfiles}% 注意这里的文件路径不能用 ./main.tex ,否则用latexmk编译子文件会报错
\graphicspath{{\subfix{./image/}}{\subfix{../../image/}}} % 指定图片引用路径,后续可以直接使用图片文件名
% 如果图片存放在上级目录,则使用 ../../image/ 作为备用路径

% 例如:
% \begin{figure}[H]
% \centering
% \includegraphics[width=0.3\textwidth]{图.png}
% \caption{}
% \label{figure:图}
% \end{figure}
% 注意:上述\label{}一定要放在\caption{}之后,否则引用图片序号会只会显示??.

\begin{document}

\section{次调和函数}

\begin{definition}
设 \(D\) 是 \(\mathbb{C}\) 中的域。如果 \(D\) 上的实值函数 \(u: D \to \mathbb{R} \cup \{-\infty\}\) (\(u \not\equiv -\infty\)) 满足
\begin{enumerate}[(1)]
\item \(u\) 是上半连续的，即对 \(D\) 中的每一点 \(a\)，有 \(\overline{\lim\limits_{z \to a}} u(z) \leqslant u(a)\)；
\item 对任意以 \(a\) 为中心、\(r\) 为半径的闭圆盘 \(\overline{B(a,r)} \subseteq D\)，有不等式
\begin{align}
u(a) \leqslant \frac{1}{2\pi} \int_{0}^{2\pi} u(a+re^{i\theta}) \mathrm{d}\theta,
\label{eq:submean}
\end{align}
\end{enumerate}
则称 \(u\) 是 \(D\) 上的\textbf{次调和函数}. 满足不等式 \eqref{eq:submean} 的函数称为具有\textbf{次平均值性质}.
\end{definition}
\begin{remark}
次调和函数是凸函数概念在平面上的推广. 事实上, 如果把 \(\frac{\mathrm{d}^2 u}{\mathrm{d}x^2} = 0\) 看作一维的 Laplace 方程, 那么这个方程的解 \(u = ax + b\) 便是一维的调和函数. 而凸函数是在任一区间的两个端点处与一线性函数有相同的值, 在区间内部, 它不超过这个线性函数. 把区间换成平面上的区域, 线性函数换成二维调和函数, 那么凸函数就是这里定义的次调和函数.
\end{remark}
\begin{remark}
每个调和函数一定是次调和的. 因此, 次调和函数是比调和函数更宽的一类函数, 但它仍具有调和函数的一些优良性质, 下面的极值特征便是其中之一.
\end{remark}

\begin{theorem}
\label{theorem:max_principle}
设 \(D\) 是 \(\mathbb{C}\) 中的域, \(u\) 是 \(D\) 上的连续实值函数. \(u\) 是 \(D\) 上的次调和函数的充分必要条件是, 对任意域 \(G \subset\subset D\) 及任意在 \(\overline{G}\) 上连续、在 \(G\) 内调和的函数 \(h\), 如果 \(u(z) \leqslant h(z)\) 在 \(\partial G\) 上成立, 那么在 \(G\) 内也有 \(u(z) \leqslant h(z)\).
\end{theorem}
\begin{proof}

{\heiti 必要性:} 若对某个 \(z_0 \in G\) 有 \(u(z_0) > h(z_0)\) 成立, 令 \(u_1 = u - h\), 则 \(u_1(z_0) > 0\). 因为 \(u_1\) 在 \(\overline{G}\) 上连续, 故在 \(\overline{G}\) 上达到它的最大值 \(M\), 记 \(E = \{ z \in \overline{G} : u_1(z) = M \}\). 由于在 \(\partial G\) 上 \(u_1 \leqslant 0\), 所以 \(u_1\) 的最大值只能在 \(G\) 内取到, 因此 \(E\) 是 \(G\) 中的紧子集. 今设 \(a\) 是 \(E\) 的一个边界点, 于是有 \(r>0\), 使得 \(\overline{B(a,r)} \subseteq G\). 但 \(\overline{B(a,r)}\) 的边界上必有某段弧不在 \(E\) 中, 因而
\begin{align}
u_1(a) = M > \frac{1}{2\pi} \int_{0}^{2\pi} u_1(a+re^{i\theta}) \mathrm{d}\theta.
\label{eq:contradict}
\end{align}
另一方面, \(u\) 和 \(h\) 分别在 \(G\) 中次调和与调和, 因而有
\begin{align*}
u(a) \leqslant \frac{1}{2\pi} \int_{0}^{2\pi} u(a+re^{i\theta}) \mathrm{d}\theta, \\
h(a) = \frac{1}{2\pi} \int_{0}^{2\pi} h(a+re^{i\theta}) \mathrm{d}\theta.
\end{align*}
由此即得
\begin{align*}
u_1(a) \leqslant \frac{1}{2\pi} \int_{0}^{2\pi} u_1(a+re^{i\theta}) \mathrm{d}\theta.
\end{align*}
这和 \eqref{eq:contradict} 式矛盾.

{\heiti 充分性:} 任取 \(\overline{B(a,r)} \subseteq D\), 那么存在 \(B(a,r)\) 中的调和函数 \(h\), 它在圆周上和 \(u\) 一致. 于是由假定, \(u(z) \leqslant h(z)\) 在圆内成立. 这样
\begin{align*}
u(a) \leqslant h(a) = \frac{1}{2\pi} \int_{0}^{2\pi} h(a+re^{i\theta}) \mathrm{d}\theta = \frac{1}{2\pi} \int_{0}^{2\pi} u(a+re^{i\theta}) \mathrm{d}\theta.
\end{align*}
这正好说明 \(u\) 是次调和函数.
\end{proof}

\begin{theorem}
\label{theorem:mean_monotonicity}
设 \(u\) 是单位圆盘 \(U\) 中的次调和函数, 令
\begin{align*}
m(r) = \frac{1}{2\pi} \int_{0}^{2\pi} u(re^{i\theta}) \mathrm{d}\theta, \quad 0 \leqslant r < 1,
\end{align*}
那么 \(m(r)\) 是 \(r\) 的非降函数.
\end{theorem}

\begin{proof}
设 \(0 < r_1 < r_2 < 1\). 存在圆盘 \(B(0,r_2)\) 上的调和函数 \(h\), 它在圆周上和 \(u\) 一致. 由 \cref{theorem:max_principle}, 在 \(B(0,r_2)\) 中有 \(u \leqslant h\), 因而
\begin{align*}
m(r_1) &= \frac{1}{2\pi} \int_{0}^{2\pi} u(r_1e^{i\theta}) \mathrm{d}\theta \leqslant \frac{1}{2\pi} \int_{0}^{2\pi} h(r_1e^{i\theta}) \mathrm{d}\theta \\
&= h(0) = \frac{1}{2\pi} \int_{0}^{2\pi} h(r_2e^{i\theta}) \mathrm{d}\theta \\
&= \frac{1}{2\pi} \int_{0}^{2\pi} u(r_2e^{i\theta}) \mathrm{d}\theta = m(r_2).
\end{align*}\end{proof}

\begin{proposition}
\label{proposition:convex_composition}
设 \(u\) 是域 \(D\) 上的次调和函数, \(\varphi\) 是 \((\infty, \infty)\) 上递增的凸函数, 那么 \(\varphi \circ u\) 也是 \(D\) 上的次调和函数.
\end{proposition}

\begin{proof}
因为 \(u\) 是次调和函数, 故对任意 \(\overline{B(a,r)} \subseteq D\), 有
\begin{align*}
u(a) \leqslant \frac{1}{2\pi} \int_{0}^{2\pi} u(a+re^{i\theta}) \mathrm{d}\theta.
\end{align*}
又因为 \(\varphi\) 是递增的凸函数, 所以
\begin{align*}
\varphi(u(a)) \leqslant \varphi\left( \frac{1}{2\pi} \int_{0}^{2\pi} u(a+re^{i\theta}) \mathrm{d}\theta \right) \leqslant \frac{1}{2\pi} \int_{0}^{2\pi} \varphi(u(a+re^{i\theta})) \mathrm{d}\theta.
\end{align*}
因而 \(\varphi \circ u\) 也是次调和函数.
\end{proof}

\begin{proposition}
\label{proposition:analytic_subharmonic}
设 \(f\) 是域 \(D\) 上的全纯函数, \(f \not\equiv 0\), 那么 \(\ln |f|\) 和 \(|f|^p\) (\(0 < p < \infty\)) 都是 \(D\) 上的次调和函数.
\end{proposition}

\begin{proof}
容易知道, \(\ln |f|\) 是上半连续的. 任取圆盘 \(\overline{B(a,r)} \subseteq D\), 我们要证明
\begin{align}
\ln |f(a)| \leqslant \frac{1}{2\pi} \int_{0}^{2\pi} \ln |f(a+re^{i\theta})| \mathrm{d}\theta.
\label{eq:log_subharmonic}
\end{align}
如果 \(f(a) = 0\), 那么 \(\ln |f(a)| = -\infty\), \eqref{eq:log_subharmonic} 式显然成立. 今设 \(f(a) \neq 0\), \(f\) 在 \(\overline{B(a,r)}\) 中也没有其他零点, 那么通过直接计算知道, \(\ln |f|\) 是 \(B(a,r)\) 中的调和函数, 因而它是次调和的. 现若 \(f\) 在 \(B(a,r)\) 中有零点 \(z_1, \cdots, z_m\), 重零点按重数重复计算, 根据 Jensen 公式, 有
\begin{align*}
\ln |f(a)| &= - \sum\limits_{k=1}^{m} \ln \frac{r}{|a - z_k|} + \frac{1}{2\pi} \int_{0}^{2\pi} \ln |f(a+re^{i\theta})| \mathrm{d}\theta \\
&\leqslant \frac{1}{2\pi} \int_{0}^{2\pi} \ln |f(a+re^{i\theta})| \mathrm{d}\theta.
\end{align*}
这就是 \eqref{eq:log_subharmonic} 式. 这就证明了 \(\ln |f|\) 的次调和性.

因为 \(\varphi(t) = e^{pt}\) 是递增的凸函数, 而 \(|f|^p = e^{p \ln |f|} = \varphi(\ln |f|)\), 故由 \cref{proposition:convex_composition}, \(|f|^p\) (\(0 < p < \infty\)) 是次调和的.
\end{proof}

\begin{proposition}
设 \(f \in H(U)\), 那么对任意 \(0 < p < \infty\), 积分平均
\begin{align*}
M_p(r,f) = \left( \frac{1}{2\pi} \int_{0}^{2\pi} |f(re^{i\theta})|^p \mathrm{d}\theta \right)^{\frac{1}{p}}
\end{align*}
是 \(r\) 的非降函数.
\end{proposition}

\begin{proposition}
\label{proposition:Delta_positive_max}
设 \(u \in C^2(D)\) 是一实值函数, 如果对任意 \(z \in D\), \(\Delta u(z) \geqslant 0\), 那么对任意域 \(G \subset\subset D\), \(u\) 在 \(G\) 上的最大值必在 \(\partial G\) 上取到.
\end{proposition}
\begin{proof}

先设对每点 \(z \in D\), 有 \(\Delta u(z) > 0\). 如果 \(u\) 在 \(G\) 上的最大值在 \(G\) 的内点 \(z_0\) 处取到, 记 \(z_0 = x_0 + i y_0\), \(g(t) = u(x_0,t)\). 那么 \(g\) 在 \(t=y_0\) 处有极大值, 因而 \(\frac{\partial^2 u}{\partial y^2}(z_0) = g''(y_0) \leqslant 0\). 同理, \(\frac{\partial^2 u}{\partial x^2}(z_0) \leqslant 0\). 所以 \(\Delta u(z_0) \leqslant 0\), 这与 \(\Delta u(z_0) > 0\) 矛盾.

现设 \(\Delta u(z) \geqslant 0\). 令 \(u_{\varepsilon}(z) = u(z) + \varepsilon |z|^2, \varepsilon > 0\). 于是, \(\Delta u_{\varepsilon}(z) = \Delta u(z) + \varepsilon > 0\). 因而 \(u_\varepsilon\) 在 \(\overline{G}\) 上的最大值必在 \(\partial G\) 上取到, 即 \(u_\varepsilon(z) \leqslant \sup\limits_{\zeta \in \partial G} u_\varepsilon(\zeta)\). 令 \(\varepsilon \to 0\), 即得 \(u(z) \leqslant \sup\limits_{\zeta \in \partial G} u(\zeta)\). 这便证明了结论.
\end{proof}

\begin{theorem}
\label{theorem:laplacian_characterization}
设 \(u \in C^2(D)\) 是一实值函数, 那么 \(u\) 是 \(D\) 上的次调和函数的充分必要条件是, 对任意 \(z \in D\), 有 \(\Delta u(z) \geqslant 0\).
\end{theorem}
\begin{proof}

{\heiti 充分性:} 设 \(\Delta u(z) \geqslant 0\) 在 \(D\) 内处处成立. 对于 \(G \subset\subset D\) 上的调和函数 \(h\), 如果 \(u \leqslant h\) 在 \(\partial G\) 上成立, 我们要证明 \(u \leqslant h\) 在 \(G\) 内成立. 令 \(u_1 = u - h\), 那么
\begin{align*}
\Delta u_1(z) = \Delta(u-h)(z) = \Delta u(z) \geqslant 0.
\end{align*}
而在 \(\partial G\) 上 \(u_1 \leqslant 0\), 于是由 \cref{proposition:Delta_positive_max}, 在 \(G\) 内也有 \(u_1 \leqslant 0\), 即 \(u \leqslant h\) 在 \(G\) 内成立. 故由 \cref{theorem:max_principle} 知, \(u\) 是 \(D\) 上的次调和函数.

{\heiti 必要性:} 设 \(u\) 是 \(D\) 上的次调和函数. 如果存在 \(a \in D\), 使得 \(\Delta u(a) < 0\), 那么有 \(a\) 的一个邻域 \(B(a,\varepsilon)\) 使 \(\Delta u(z) < 0\) 对任意 \(z \in B(a,\varepsilon)\) 成立, 即 \(-\Delta u(z) > 0\) 在 \(B(a,\varepsilon)\) 中成立. 由充分性证明的结果, \(-u\) 是 \(B(a,\varepsilon)\) 上的次调和函数, 因而 \(u\) 的平均值公式在 \(B(a,\varepsilon)\) 中的任意小圆盘上成立. 由调和函数的性质知, \(u\) 是 \(B(a,\varepsilon)\) 中的调和函数, 从而 \(\Delta u(a) = 0\), 这与假设 \(\Delta u(a) < 0\) 矛盾. 所以, 对任意 \(z \in D\), \(\Delta u(z) \geqslant 0\) 处处成立.
\end{proof}



\end{document}